\documentclass[11pt,a4paper]{article}
\usepackage{cursus}
\author{dr. Maarten Slembrouck}

\begin{document}

\labotitel{Labo 4 -- Dataverwerking II}{Combineren, samenvatten, herschikken}{Cursustekst bij het notebook \code{lab4\_data\_processing\_II.py}}

\section*{Waarover gaat dit labo?}

Vorige week maakte je van rommelige exports nette tabellen: één rij per
atleet, één rij per test. Maar de vragen van een coach gaan zelden over
één rij. ``Hoe groot zijn onze verdedigers gemiddeld?'' gaat over een
\emph{groep} rijen. ``Hoe snel sprinten onze keepers?'' heeft informatie
uit \emph{twee} tabellen nodig. ``Is iedereen vooruitgegaan sinds de
eerste test?'' vraagt om de tabel anders te \emph{schikken}. En de
hartslagdata van een training is een \emph{tijdreeks}, met één meting per
seconde.

\textbf{Na dit labo kun je} data per groep samenvatten met
\code{groupby}, tabellen koppelen met \code{merge}, data herschikken met
\code{melt} en \code{pivot\_table}, en eenvoudige kenmerken van een
tijdreeks berekenen (hartslagzones, voortschrijdend gemiddelde,
hersampling).

Net als vorige week is elk voorbeeld echt uitgevoerd op de data van de
club. Lees deze tekst naast het notebook en typ de voorbeelden zelf over.

% ---------------------------------------------------------------------------
\section{Vertrekpunt: de opgekuiste data van vorige week}

Het notebook begint met de opkuis van labo 3, in compacte vorm. Twee
hulpfuncties doen het werk: \code{clean\_positions} vertaalt de 13
schrijfwijzen naar 4 posities, \code{fix\_sprint\_units} zet
milliseconden om naar seconden.

\begin{python}[noout]
import pandas as pd

def clean_positions(position_column):
    position_map = {
        "forward": "Forward", "fwd": "Forward",
        "midfielder": "Midfielder", "mid": "Midfielder",
        "defender": "Defender", "def": "Defender",
        "goalkeeper": "Goalkeeper", "gk": "Goalkeeper",
    }
    return position_column.str.strip().str.lower().map(position_map)

def fix_sprint_units(seconds_column):
    return seconds_column.where(seconds_column < 20, seconds_column / 1000)

REFERENCE_DATE = pd.Timestamp("2026-09-14")

athletes = pd.read_csv(f"{DATA_DIR}/athletes.csv").assign(
    position=lambda d: clean_positions(d["position"]),
    weight_kg=lambda d: d["weight_kg"].str.replace(",", ".").astype(float),
    age=lambda d: (REFERENCE_DATE - pd.to_datetime(d["birthdate"])).dt.days // 365,
)
tests = pd.read_csv(f"{DATA_DIR}/tests.csv").assign(
    sprint_10m_s=lambda d: fix_sprint_units(d["sprint_10m_s"]),
    sprint_30m_s=lambda d: fix_sprint_units(d["sprint_30m_s"]),
)
\end{python}

Lees deze code rustig: alles erin heb je vorige week gezien. Merk op dat
de dubbele atleet hier \emph{bewust} nog niet verwijderd is. Dat wordt
belangrijk in sectie~\ref{sec:merge}.

\subsection*{Herhaling: foutmeldingen lezen}

Het notebook start met een kapotte cel die je moet herstellen. Een
foutmelding ziet er eng uit, maar de belangrijkste informatie staat
altijd op de \textbf{laatste regel}: het soort fout en een korte uitleg.
Twee fouten die je vandaag zeker tegenkomt:

\begin{python}[error]
athletes["positio"]
\end{python}

Een \code{KeyError} betekent: deze sleutel (hier: kolomnaam) bestaat
niet. Bijna altijd een tikfout.

\begin{python}[error]
pd.Series(["57,5", "60,3"]).astype(float)
\end{python}

Een \code{ValueError} bij \code{astype(float)} betekent: deze tekst kan
ik niet omzetten naar een getal. Hier staat de komma in de weg.

% ---------------------------------------------------------------------------
\section{Groepen samenvatten met \code{groupby}}

\subsection{Splitsen, toepassen, combineren}

Stel dat de coach de gemiddelde leeftijd per positie wil kennen. In
gedachten doe je drie dingen:

\begin{enumerate}
  \item \textbf{splitsen}: de atleten in vier groepjes verdelen, één per
        positie;
  \item \textbf{toepassen}: per groepje het gemiddelde berekenen;
  \item \textbf{combineren}: de vier gemiddelden in één tabel zetten.
\end{enumerate}

\begin{center}
\begin{tikzpicture}[font=\small\sffamily, box/.style={rounded corners=2pt, minimum width=2.3cm, minimum height=0.6cm, text=white}]
  \node[draw=none, fill=tint, rounded corners=3pt, minimum width=2.3cm, minimum height=3.2cm, text=pitch, align=center] (all) at (0,0) {alle\\atleten};
  \foreach \p/\c/\y in {Forward/forward/1.2, Midfielder/midfielder/0.4, Defender/defender/-0.4, Goalkeeper/goalkeeper/-1.2} {
    \node[box, fill=\c] (\p) at (3.6,\y) {\p};
    \draw[-{Stealth}, muted] (all.east) -- (\p.west);
    \node[text=muted, font=\small\itshape] (m\p) at (6.1,\y) {gemiddelde};
    \draw[-{Stealth}, muted] (\p.east) -- (m\p.west);
    \draw[-{Stealth}, muted] (m\p.east) -- (8.2,0);
  }
  \node[draw=none, fill=pitch, rounded corners=3pt, minimum width=2.1cm, minimum height=2.0cm, text=white, align=center] at (9.3,0) {één rij\\per positie};
  \node[text=pitch, font=\small\bfseries\sffamily] at (1.8,2.1) {splitsen};
  \node[text=pitch, font=\small\bfseries\sffamily] at (5.0,2.1) {toepassen};
  \node[text=pitch, font=\small\bfseries\sffamily] at (8.4,2.1) {combineren};
\end{tikzpicture}
\end{center}

In pandas is dat één regel. \code{groupby("position")} splitst,
\code{["age"]} kiest de kolom, \code{.mean()} past het gemiddelde toe en
combineert:

\begin{python}
mean_age_per_position = athletes.groupby("position")["age"].mean().round(1)
mean_age_per_position
\end{python}

Het resultaat is een Series met de posities als \emph{index}. In plaats
van \code{.mean()} kun je ook \code{.median()}, \code{.min()},
\code{.max()}, \code{.sum()} of \code{.count()} gebruiken.

\begin{opgelet}[\code{size()} of \code{count()}?]
\code{.size()} telt het aantal rijen per groep. \code{.count()} telt het
aantal \emph{ingevulde} waarden in een kolom en slaat \code{NaN} over.
Wil je weten hoeveel atleten er per positie zijn, gebruik dan
\code{.size()}.
\end{opgelet}

\begin{python}
athletes.groupby("position").size()
\end{python}

\subsection{Getallen in categorieën: \code{pd.cut}}

Jeugdploegen worden ingedeeld in leeftijdscategorieën: U13, U15, U17,
U19. \code{pd.cut} verdeelt een getallenkolom in vakjes. Je geeft de
grenzen (\code{bins}) en de namen van de vakjes (\code{labels}) mee:

\begin{python}
athletes_with_age_group = athletes.assign(
    age_group=lambda d: pd.cut(
        d["age"], bins=[10, 12, 14, 16, 18], labels=["U13", "U15", "U17", "U19"]
    )
)
athletes_with_age_group[["name", "age", "age_group"]].head()
\end{python}

Vijf grenzen maken vier vakjes: van 10 tot 12, van 12 tot 14, enzovoort.
Een vakje bevat zijn \emph{rechtergrens} maar niet zijn linkergrens: een
atleet van 12 valt in U13 (10 tot en met 12), een atleet van 13 in U15.
Er is dus altijd één label minder dan er grenzen zijn.

\subsection{Groeperen op meerdere kolommen en \code{agg}}

Je kunt op twee kolommen tegelijk groeperen door een lijst mee te geven.
Met \code{.agg(...)} bereken je meerdere statistieken in één keer en geef
je elke nieuwe kolom zelf een naam:

\begin{python}
athletes_with_age_group.groupby(["position", "age_group"], observed=True).agg(
    n_athletes=("athlete_id", "size"),
    mean_height_cm=("height_cm", "mean"),
)
\end{python}

Lees \code{mean\_height\_cm=("height\_cm", "mean")} als: ``maak een
kolom \code{mean\_height\_cm}, door op de kolom \code{height\_cm} de
functie \code{mean} toe te passen''. Een paar extra dingen:

\begin{itemize}
  \item De index heeft nu twee niveaus (positie en leeftijdscategorie):
        een \emph{MultiIndex}. Met \code{.reset\_index()} maak je er weer
        gewone kolommen van.
  \item \code{age\_group} is een \emph{categorische} kolom (gemaakt door
        \code{pd.cut}). \code{observed=True} zorgt ervoor dat combinaties
        zonder atleten (bv.\ keepers in U13) niet als lege rij verschijnen.
\end{itemize}

\begin{probeer}
Bereken de maximale lengte per positie. Groepeer daarna op geslacht en
tel met \code{.size()} hoeveel jongens en meisjes er zijn.
\end{probeer}

% ---------------------------------------------------------------------------
\section{Tabellen koppelen met \code{merge}}\label{sec:merge}

\subsection{Een sleutel in twee tabellen}

\code{athletes} weet van elke atleet de positie, \code{tests} kent de
sprinttijden. Om de sprinttijd per positie te bekijken, moet je de twee
koppelen. Dat kan omdat beide tabellen een kolom \code{athlete\_id}
hebben: de \textbf{sleutel}.

Een klein voorbeeld maakt duidelijk wat er gebeurt:

\begin{python}
people = pd.DataFrame({"athlete_id": ["A1", "A2", "A3"],
                       "name": ["Yana", "Stan", "Emma"]})
results = pd.DataFrame({"athlete_id": ["A1", "A1", "A2", "A9"],
                        "sprint_s": [5.3, 5.2, 5.9, 6.1]})
people.merge(results, on="athlete_id", how="inner")
\end{python}

Voor elke rij in \code{people} zoekt pandas de rijen in \code{results}
met dezelfde \code{athlete\_id}. Yana heeft twee resultaten, dus ze
verschijnt twee keer. Emma (A3) heeft geen resultaat en A9 heeft geen
naam: met \code{how="inner"} vallen ze allebei weg.

\begin{python}
people.merge(results, on="athlete_id", how="left")
\end{python}

Met \code{how="left"} blijft \emph{elke} rij van de linkertabel
behouden. Emma staat erin, met \code{NaN} als sprinttijd.

\begin{kernidee}[De soorten joins]
\textbf{inner}: alleen sleutels die in beide tabellen staan.
\textbf{left}: alle rijen van de linkertabel, aangevuld waar mogelijk.
\textbf{right}: hetzelfde, vanuit de rechtertabel.
\textbf{outer}: alles uit beide tabellen.
Controleer na elke merge het aantal rijen: dat vertelt je of er iets
wegviel of verdubbelde.
\end{kernidee}

Met de echte data:

\begin{python}
merged_inner = athletes.merge(tests, on="athlete_id", how="inner")
merged_left = athletes.merge(tests, on="athlete_id", how="left")
(len(athletes), len(tests), len(merged_inner), len(merged_left))
\end{python}

Hier zijn beide even lang, omdat elke atleet minstens één testresultaat
heeft. Met echte data is dat zelden gegarandeerd.

\subsection{De dubbele atleet}

Eén atleet staat twee keer in \code{athletes}, met twee verschillende
ID's. Je vindt die rijen met \code{duplicated}; met \code{keep=False}
worden \emph{alle} exemplaren gemarkeerd, niet alleen de tweede:

\begin{python}
same_person = athletes.duplicated(subset=["name", "birthdate"], keep=False)
athletes.loc[same_person, ["athlete_id", "name", "birthdate"]]
\end{python}

De merge doet perfect wat gevraagd is, maar weet niet dat het om dezelfde
persoon gaat. Na de merge staan de testen van die atleet onder twee ID's:

\begin{python}
dupe_name = athletes.loc[same_person, "name"].iloc[0]
merged_inner.loc[merged_inner["name"] == dupe_name,
                 ["athlete_id", "name", "test_moment", "sprint_30m_s"]]
\end{python}

Elke statistiek die je na deze merge berekent, telt deze atleet dubbel.
De oplossing: verwijder de dubbels \emph{vóór} de merge, op naam en
geboortedatum. Na de merge kan dat niet meer, want dan delen ook de
verschillende testmomenten van één atleet dezelfde naam en
geboortedatum: je zou echte testresultaten weggooien.

\begin{python}
athletes_unique = athletes.drop_duplicates(subset=["name", "birthdate"], keep="first")
combined = athletes_unique.merge(tests, on="athlete_id", how="inner")
combined.shape
\end{python}

\subsection{Verzamelingen en \code{isin}}

Om te weten welke rijen bij het ontdubbelen verdwenen, heb je twee
hulpmiddelen nodig.

Een \textbf{verzameling} (\code{set}) is een lijst zonder dubbels en
zonder volgorde. Het verschil van twee verzamelingen (\code{a - b}) geeft
de elementen die in \code{a} zitten maar niet in \code{b}:

\begin{python}
set(["A1", "A2", "A3"]) - set(["A2"])
\end{python}

\code{.isin(...)} geeft een masker dat \code{True} is voor elke rij
waarvan de waarde in de opgegeven lijst of verzameling staat:

\begin{python}
results["athlete_id"].isin({"A1", "A9"})
\end{python}

Met deze twee bouwstenen (en \code{.sum()} op een masker) kun je
oefening 2 in het notebook oplossen.

% ---------------------------------------------------------------------------
\section{Lang en breed: \code{pivot\_table} en \code{melt}}

\subsection{Twee vormen voor dezelfde data}

\code{combined} is \textbf{lang}: één rij per atleet \emph{per
testmoment}. Een atleet met drie testen staat er drie keer in. Dezelfde
informatie kan ook \textbf{breed}: één rij per atleet, met een kolom per
testmoment.

\begin{python}
sprint_wide = combined.pivot_table(
    index="athlete_id", columns="test_moment", values="sprint_30m_s"
)
sprint_wide.head()
\end{python}

De drie argumenten lees je zo:
\begin{itemize}
  \item \code{index}: wat wordt een rij (één per atleet);
  \item \code{columns}: welke waarden worden kolommen (de testmomenten
        1, 2 en 3);
  \item \code{values}: welke waarde komt in de cellen (de sprinttijd).
\end{itemize}

Ontbreekt een combinatie, dan krijg je \code{NaN}. Zijn er meerdere rijen
voor dezelfde combinatie, dan neemt \code{pivot\_table} standaard het
gemiddelde.

\subsection{Waarom breed handig is}

In de brede vorm is vooruitgang tussen twee testmomenten één
aftrekking. Let op: de kolomnamen zijn hier \emph{getallen}, geen tekst.
Je schrijft dus \code{sprint\_wide[2]}, niet \code{sprint\_wide["2"]}.

\begin{python}
sprint_progress_1_to_2 = sprint_wide[2] - sprint_wide[1]
sprint_progress_1_to_2.describe()
\end{python}

Een negatief getal betekent hier vooruitgang: een lagere sprinttijd is
sneller. Bij spronghoogte is het omgekeerd: daar is hoger beter.

\subsection{Terug naar lang: \code{melt}}

\code{melt} doet het omgekeerde van \code{pivot\_table}: kolommen worden
weer rijen. Plotly express (labo 5) werkt het liefst met lange data.
Eerst maak je met \code{.reset\_index()} van de index weer een gewone
kolom:

\begin{python}
sprint_long_again = sprint_wide.reset_index().melt(
    id_vars="athlete_id", var_name="test_moment", value_name="sprint_30m_s"
)
sprint_long_again.head()
\end{python}

\code{id\_vars} zijn de kolommen die blijven staan, \code{var\_name} is
de naam van de nieuwe kolom met de vroegere kolomnamen, \code{value\_name}
de naam van de kolom met de waarden.

% ---------------------------------------------------------------------------
\section{Tijdreeksen: de hartslag tijdens een training}

\code{hr\_session.csv} bevat per seconde de hartslag en snelheid van vijf
atleten tijdens één training. Een tabel waarin elke rij een moment in de
tijd is, heet een \textbf{tijdreeks}.

\subsection{Een tijdstempel maken}

Het bestand heeft een datum en het aantal seconden sinds de start
(\code{time\_s}). Samen maken ze een tijdstempel.
\code{pd.to\_timedelta(..., unit="s")} zet een aantal seconden om in een
tijdsduur, die je bij een datum kunt optellen:

\begin{python}
hr_session = pd.read_csv(f"{DATA_DIR}/hr_session.csv")
hr_session["timestamp"] = pd.to_datetime(hr_session["date"]) + pd.to_timedelta(
    hr_session["time_s"], unit="s"
)
hr_session.head(3)
\end{python}

\subsection{Hartslagzones}

Trainers drukken intensiteit uit als percentage van de \emph{maximale}
hartslag. Een veelgebruikte vuistregel schat die als 220 min de leeftijd.
De leeftijd staat in \code{athletes}, dus eerst een merge:

\begin{python}
ages = athletes[["athlete_id", "age"]]
hr_with_age = hr_session.merge(ages, on="athlete_id", how="left")
hr_with_age["hr_max"] = 220 - hr_with_age["age"]
hr_with_age["pct_hr_max"] = (
    100 * hr_with_age["heart_rate_bpm"] / hr_with_age["hr_max"]
)
hr_with_age[["athlete_id", "heart_rate_bpm", "hr_max", "pct_hr_max"]].head(3)
\end{python}

Met \code{pd.cut} deel je het percentage in vijf zones in. Het merkwaardige
getal 1000 als laatste grens zorgt ervoor dat ook een percentage boven
100 nog in zone 5 valt:

\begin{python}
hr_zoned = hr_with_age.assign(
    hr_zone=lambda d: pd.cut(
        d["pct_hr_max"],
        bins=[0, 60, 70, 80, 90, 1000],
        labels=["Zone 1", "Zone 2", "Zone 3", "Zone 4", "Zone 5"],
    )
)
hr_zoned["hr_zone"].value_counts()
\end{python}

\subsection{Tijd per zone}

Elke rij is één seconde. Het aantal rijen per atleet en per zone, gedeeld
door 60, geeft dus het aantal minuten:

\begin{python}
time_per_zone_minutes = (
    hr_zoned.groupby(["athlete_id", "hr_zone"], observed=True).size() / 60
).rename("minutes").reset_index()
time_per_zone_minutes.head(5)
\end{python}

Wil je enkel bepaalde zones, filter dan eerst de rijen met
\code{.isin([...])} en groepeer daarna. Dat is de kern van oefening 4.

\subsection{Voortschrijdend gemiddelde en hersampling}

Een hartslag per seconde springt op en neer. Twee technieken maken de
reeks rustiger. Beide hebben een tijdreeks nodig die \emph{gesorteerd} is
op tijd, met de tijd als index (\code{set\_index("timestamp")}):

\begin{python}[noout]
one_athlete = hr_zoned.loc[hr_zoned["athlete_id"] == "A001"]
one_athlete = one_athlete.sort_values("timestamp")
hr_series = one_athlete.set_index("timestamp")["heart_rate_bpm"]
\end{python}

Een \textbf{voortschrijdend gemiddelde} (\code{rolling}) berekent op elk
moment het gemiddelde van de voorbije 30 seconden. Je houdt even veel
rijen, maar de lijn wordt vloeiender:

\begin{python}
hr_series.rolling("30s").mean().iloc[[0, 1, 29, 30, 31]]
\end{python}

\textbf{Hersampling} (\code{resample}) groepeert de tijdreeks in vaste
blokken, hier van één minuut, en vat elk blok samen. Je krijgt veel
minder rijen: één per minuut.

\begin{python}
hr_series.resample("1min").mean().head(4)
\end{python}

\begin{kernidee}[\code{rolling} of \code{resample}?]
\code{rolling} geeft een waarde per oorspronkelijke rij (een glijdend
venster); handig om een grafiek vloeiender te maken.
\code{resample} geeft een waarde per tijdsblok; handig om per minuut of
per kwartier samen te vatten. Je kunt \code{resample} zien als een
\code{groupby} op tijd.
\end{kernidee}

% ---------------------------------------------------------------------------
\section{De minicase: een trainingsrapport}

In de minicase (per twee) maak je één rij per atleet met de minuten in
zone 4--5, de maximale snelheid en de afgelegde afstand. Twee ideeën
helpen je daarbij.

\textbf{Afstand uit snelheid.} De snelheid staat in km/u en er is één
meting per seconde. In één seconde leg je dus \code{speed\_kmh / 3600}
kilometer af. De som van die afstanden over de hele training is de totale
afstand.

\textbf{Series combineren via de index.} Als twee Series dezelfde index
hebben (bv.\ de \code{athlete\_id}), plaatst pandas de waarden bij een
toewijzing of berekening automatisch naast de juiste rij, ook al is de
volgorde anders. Dat heet \emph{uitlijnen op de index}:

\begin{python}
report = pd.DataFrame({"max_speed": [27.1, 24.5]}, index=["A002", "A001"])
extra = pd.Series([10.0, 12.5], index=["A001", "A002"])
report["minutes"] = extra
report
\end{python}

A001 krijgt 10.0 en A002 krijgt 12.5, ook al stonden ze in
\code{extra} in een andere volgorde dan in \code{report}.

% ---------------------------------------------------------------------------
\section{Samenvatting}

\begin{center}
\rowcolors{2}{tint}{white}
\begin{tabularx}{\textwidth}{@{}L{8.4cm}Y@{}}
\toprule
\textbf{Code} & \textbf{Wat het doet} \\
\midrule
\code{df.groupby("a")["b"].mean()} & gemiddelde van \code{b} per waarde van \code{a} \\
\code{df.groupby("a").size()} & aantal rijen per groep \\
\code{df.groupby(["a", "b"]).agg(naam=("kol", "mean"))} & meerdere benoemde statistieken per groep \\
\code{pd.cut(kol, bins=..., labels=...)} & getallen indelen in categorieën \\
\code{.reset\_index()} & index terug omzetten naar gewone kolommen \\
\code{links.merge(rechts, on="sleutel", how="inner")} & twee tabellen koppelen op een sleutel \\
\code{set(a) - set(b)}, \code{kol.isin(lijst)} & verschil van verzamelingen, masker ``zit in'' \\
\code{df.pivot\_table(index=, columns=, values=)} & van lang naar breed \\
\code{df.melt(id\_vars=, var\_name=, value\_name=)} & van breed naar lang \\
\code{pd.to\_timedelta(kol, unit="s")} & aantal seconden naar tijdsduur \\
\code{s.rolling("30s").mean()} & voortschrijdend gemiddelde over 30 seconden \\
\code{s.resample("1min").mean()} & gemiddelde per blok van één minuut \\
\bottomrule
\end{tabularx}
\end{center}

\subsection*{Veelgemaakte fouten}

\begin{itemize}
  \item Na een merge het aantal rijen niet controleren.
  \item Dubbels verwijderen \emph{na} de merge in plaats van ervoor.
  \item \code{.count()} gebruiken om rijen te tellen terwijl er
        ontbrekende waarden zijn; gebruik \code{.size()}.
  \item Kolomnamen na een \code{pivot\_table} als tekst schrijven
        (\code{wide["2"]}) terwijl het getallen zijn (\code{wide[2]}).
  \item \code{rolling} of \code{resample} gebruiken op data die niet op
        tijd gesorteerd is, of zonder tijdsindex.
\end{itemize}

\subsection*{Begrippen}

\begin{description}[leftmargin=3.2cm, labelwidth=3.0cm, font=\normalfont\bfseries]
  \item[groupby] splitsen in groepen, per groep iets berekenen, en de resultaten combineren.
  \item[aggregatie] een groep waarden samenvatten tot één waarde (gemiddelde, som, aantal).
  \item[MultiIndex] index met meerdere niveaus, bv.\ positie en leeftijdscategorie.
  \item[sleutel] kolom waarmee je twee tabellen koppelt.
  \item[join] manier van koppelen: inner, left, right of outer.
  \item[lang / breed] één rij per meting, of één rij per atleet met een kolom per meting.
  \item[tijdreeks] tabel waarin elke rij een moment in de tijd is.
\end{description}

\end{document}
